\section{Future Directions}\label{sec:chap4 future_directions}

The natural next step is to use this stratification to run stratified sampling to estimate a plan-level observable of interest $f$, for instance, democratic vote shares.
Let $\pi$ be the target distribution $e^{-J(x)}$ and $\{\phi_s\}_{s\in[S]}$ be the POU that we construct using the pipeline introduced in this paper.
We define the tilted stratum distributions by 
    \begin{align}
        \pi_s(x)\propto e^{-J_s(x)},\quad\quad J_s(x):=J(x)-\phi_s(x).
    \end{align}
Using the same method (e.g. cycle walk), we sample $S$ chains with the biased energy $J_s(x)$.
The Metropolis-Hastings steps need to be adjusted to satisfy detailed balance with respect to the tilted measures $\pi(s)$.
We then apply standard fix-point iterations methods to recover the normalizing factors $z_s$, so that we can estimate $f$ by
    \begin{align}
        \mathbb{E}_\pi[f]:=\sum_{s=1}^Sz_s\mathbb{E}_{\pi_s}[f].
    \end{align}
Ideally, each chain focus on plans around a different word $w_s\in\mathcal{S}$, which allows for more accurate estimates of the observables.

Another interesting direction is to adaptively update the alphabet as we obtain new data, so that the alphabet still captures the motifs in district space.
This can be useful when the original sample used to construct the alphabet is too small, or when the target distribution changes drastically, causing new motifs to appear.
Re-running hierarchical clustering from scratch after each new sample can be costly.
On the other hand, most hierarchical procedures are not designed for streaming data.
Online hierarchical clustering has been studied, but under different objectives~\cite{menon2019online}.
If we expect localized updates, i.e. $K$ does not change.
We can assign new data to closest centroids (i.e. Rocchio classification) and run a few $K$-median iterations (if necessary) to ensure that the letters reflect their corresponding clusters.
While this approach is intuitive and computationally cheap, the hierarchical structure is broken and it is not clear whether the theoretical guarantees~\cite{DBLP:journals/corr/abs-2409-01010} are preserved.

Finally, sampling parameters can substantially change district distributions.
For a target such as $\gamma=1.0$, we may learn better letters from an intermediate distribution, such as $\gamma=0.5$, than from a distant baseline such as $\gamma=0.0$.
We can also form a shared alphabet by taking, and possibly pruning, the union of letters learned across several values of $\gamma$.
